\documentclass[10pt,a4paper]{amsart}
\usepackage[margin=19mm]{geometry}
\usepackage{amsmath,amssymb,mathtools}
\usepackage[scr=rsfs,cal=euler]{mathalfa}
\usepackage{xfrac}
\usepackage[unicode,hidelinks,bookmarks=false]{hyperref}
\newcommand{\ie}{i.e.\ }
\newcommand{\cf}{cf.\ }
\newcommand{\resp}{respectively\ }
\newcommand{\Subsection}{\subsection}
\providecommand{\nobreakdash}{\nobreak\hbox{-}}
\setlength{\emergencystretch}{2em}
\multlinegap0pt
\usepackage{bm}
\usepackage{bbm}
\usepackage{dsfont}
\usepackage{xspace}
\usepackage{cancel}
\usepackage{mathtools}
\usepackage{amsthm}
\makeatletter
\@ifundefined{theorem}{\newtheorem{theorem}{Theorem}}{}
\@ifundefined{proposition}{\newtheorem{proposition}[theorem]{Proposition}}{}
\makeatother
\DeclareMathOperator{\GL}{GL}
\newcommand{\OK}{\mathcal{O}_K}
\title[Effective module lattices and their shortest vectors]{Effective module lattices and their shortest vectors (Definition 7)}
\author{Nihar Gargava, Vlad Serban, Maryna Viazovska, Ilaria Viglino}
\date{Source: arXiv:2402.10305v2 (CC BY 4.0)}
\begin{document}
\maketitle

\noindent\emph{Source and licence.} Effective module lattices and their shortest vectors. Version arXiv:2402.10305v2 (CC BY 4.0), distributed under the Creative Commons Attribution 4.0 International licence, as recorded in the arXiv licence metadata for this exact version and shown on the abstract page https://arxiv.org/abs/2402.10305v2. Full rights notice: licenses/arxiv-2402.10305-CC-BY-4.0.txt.\par\smallskip

\noindent\textit{Editorial scope.} Counted unit: the Proposition printed as Definition 7 of the article (source0.tex lines 289--292, source label \texttt{pr:trans\_action}), delivered together with the article's own complete original proof of it (source0.tex lines 293--330, one \texttt{proof} environment of 38 lines). No mathematics is written, repaired or re-derived: statement and proof are the article's own text line for line. Required context retained uncounted: none. Every definition, display and label of those lines is retained unchanged. Omitted from this package and not redistributed: 1--288, 331--1822. Nothing inside a retained excerpt is omitted or altered: every retained body is delivered exactly as the article's lines are written, so no omission receipt of any kind accompanies this package. Citations: the retained excerpts carry no citation command at all, so no bibliography entry of the article is transcribed and no reference is redistributed. Reference adaptation: none was needed and none was made; every reference inside the delivered unit resolves inside it, so no unresolved reference or citation is produced. Printed slips kept literal: no printed word, symbol, hypothesis or number of the counted statement or of its proof was altered. Layout and class: the article's own class is \texttt{amsart} with packages including amsfonts, pgfplots, amsmath, calc, amsthm, amssymb, hyperref, autonum, inputenc, relsize, fontenc, soul, tikz, tikz-cd; the wrapper is the lane's pinned \texttt{amsart} editorial wrapper at 10pt on A4, which restates the theorem-like environments the retained text needs and adds the article's own macro definitions that the retained text needs (\texttt{\textbackslash GL}, \texttt{\textbackslash OK}), with extra wrapper packages (bm, bbm, dsfont, xspace, cancel, mathtools). The delivered numbering is the unit's own. Assets: no retained span contains a figure environment, a graphics-inclusion command, a PDF-page inclusion, a table or a TikZ picture; no image file is copied into this package, the package declares no asset dependency and no image file is redistributed with it. Licence: version arXiv:2402.10305v2 of this article is distributed under the Creative Commons Attribution 4.0 International licence, as recorded in the arXiv licence metadata for this exact version and shown on the abstract page https://arxiv.org/abs/2402.10305v2. A full rights notice is delivered at licenses/arxiv-2402.10305-CC-BY-4.0.txt. Characters: every retained excerpt is pure ASCII, so the delivered engine, XeLaTeX with its default Latin Modern text font, reports no missing character.
\begin{proposition}
\label{pr:trans_action}
The group $\GL_{t}(K_\mathbb{R})/\mathbb{R}^{\times}$ acts transitively on the set of $\OK$-modules in the same Steinitz class. 
\end{proposition}

\begin{proof} 

	Let $\Lambda = (g , M)$ be an $\OK$-module. 

	We know from the structure theorem of finitely generated torsion-free $\OK$-modules in $K^{n}$
	(or more generally, any Dedekind domain) that there exists some pseudobasis for $M$. That is,
	for some $v_1,\dots,v_t \in M $, and some fractional ideal $\mathcal{I}\subseteq K$,
	we get 
	\begin{equation}
	  M = \OK v_1 + \OK v_2 + \dots + \OK v_{t-1} + \mathcal{I} v_t.
	\end{equation}
	This tells us that there is matrix $g \in \GL_{t}(K)$ that sends $\OK^{t-1} \times \mathcal{I}\subseteq K^{t}$ 
	to $M$. Since $\GL_{t}(K)$ preserves the Steinitz class, 
	we obtain that $[\mathcal{I}] = [\Lambda]$.

	Now if we have that for some $\OK$-modules $(g_1, M_1)$ and $(g_2,M_2)$ have the same Steinitz class, then they are both $\GL_{t}(K)$ equivalent to some 
	$(g'_1, \OK^{t-1} \times \mathcal{I})$ and
	$(g'_2, \OK^{t-1} \times \mathcal{I})$ respectively. 
	This tells us that $g'_1 ( g'_2 )^{-1} \mathbb{R}^{\times} \in \GL_{t}(K \otimes \mathbb{R})/\mathbb{R}^{\times}$ is the required map sending one to another. 

	% Let $[\Lambda] \in \cl(K)$ be the Steinitz class of a some $\OK$-module $\Lambda$.
	% It is sufficient to prove that $\GL_{t}(K)$ acts transitively on the rational $\OK$-modules $\Lambda$. For this, it is sufficent to show that for some $g \in \GL_{t}(K)$, $g\Lambda = \OK^{t-1} \times \mathcal{I}$ for some fractional ideal $\mathcal{I}$ such that $[\mathcal{I}] = [\Lambda]$. 
	
	% Let $v_1,\dots,v_t \in K^{t}$ be any $K$-linearly independent collection of vectors in $\Lambda$. Now let $g \in \GL_{t}(K)$ be a matrix whose columns are $v_1,\dots,v_t$. Then $g$ maps the vectors of $\OK^{t}$ to a sublattice inside $\Lambda$.
	% We can then define an $\OK$-ideal $\mathcal{I}$ as the largest ideal define an ideal $\mathcal{I}$ as
	% \begin{equation}
	%   \{a \in K \mid 
% g \cdot
% \begin{bmatrix}
	% \alpha_1\\
	% \vdots \\
	% \alpha_{t-1}\\
	% a
% \end{bmatrix} \in \Lambda \text{ for all }\alpha_1,\dots,\alpha_{t-1} \in \OK
	%   \}
	% \end{equation}
	
\end{proof}

\end{document}
